% !PRINTED 642
\weberpage{642}
%<<CONT
be satisfied when all of its coefficients vanish. This leads to the equations
\[
x_{\frac12\mu}=0,\ x_{\frac12\mu-1}=x_{-\frac12\mu+1},\ \dots,\ x_1=x_{-1},
\]
and every integer $\omega$ of the field $H_m$ is therefore contained in the form:
\[
\omega = x_0+x_1(r+r^{-1})+x_2(r^2+r^{-2})+\cdots
+x_{\frac12\mu-1}(r^{\frac12\mu-1}+r^{-\frac12\mu+1}).
\]

Instead of the basis (1) one may also, if one sets
\[
\varrho = r+r^{-1}
\]
choose the powers of $\varrho$:
\weq{3}{1,\ \varrho,\ \varrho^2,\ \dots,\ \varrho^{\frac12\mu-1}}
as a basis of the integers of $H_m$. For the quantities (3) can be expressed with integral coefficients in terms of (1), and conversely.

The number $\varrho$ is the root of an irreducible equation of degree $\frac12\mu$, which we shall denote by $\psi(\varrho)=0$. If $f(r)=0$ is the equation for $r$ (\S.\,169), then there subsists the identical relation
\weq{4}{f(x) = x^{\frac12\mu}\,\psi\Bigl(x+\frac1x\Bigr),}
from which, by formation of the derivative,
\weq{5}{f'(r) = r^{\frac12\mu}\,\psi'(\varrho)\,(1-r^{-2}).}

From this we can form the fundamental number $\Delta_1$ of the field $H_m$, which, since $H_m$ contains only real numbers, must be positive. This fundamental number, if we denote the norm with respect to $H_m$ by $N_1$, is:
\[
\Delta_1 = \pm\,N_1\,\psi'(\varrho).
\]

But if $N$ is the norm formed with respect to $\Omega_m$, then
\[
N\,\psi'(\varrho) = [N_1\,\psi'(\varrho)]^2 = \Delta_1^2,
\]
and accordingly formula (5), having regard to \S.\,169, yields:
\[
\Delta = N(1-r^{-2})\,\Delta_1^2.
\]

By \S.\,169, (7) and (10) we have:
\weq{6}{
\begin{array}{rll}
N(1-r^{-2}) &{}=q & \text{for odd } q\\
&{}=4 & \text{for } q=2,\\[4pt]
\pm\,\Delta &{}=q\,\Delta_1^2 & \text{for odd } q\\
&{}=4\,\Delta_1^2 & \text{for } q=2.
\end{array}
}
